% Generated draft: requires mathematical and source review.
\documentclass[11pt]{article}
\usepackage[margin=25mm]{geometry}
\usepackage{fontspec}
\setmainfont{Latin Modern Roman}
\usepackage{amsmath,amssymb,mathtools,unicode-math}
\setmathfont{Latin Modern Math}
\usepackage{enumitem,needspace,xurl,hyperref}
\setlength{\parindent}{0pt}
\setlength{\parskip}{5pt}
\setlength{\emergencystretch}{3em}
\newcommand{\N}{\mathbb{N}}
\newcommand{\Z}{\mathbb{Z}}
\newcommand{\Q}{\mathbb{Q}}
\newcommand{\R}{\mathbb{R}}
\newcommand{\C}{\mathbb{C}}
\newcommand{\F}{\mathbb{F}}
\newcommand{\A}{\mathbb{A}}
\newcommand{\PP}{\mathbb{P}}
\newcommand{\E}{\mathbb{E}}
\newcommand{\eps}{\varepsilon}
\newcommand{\dd}{\,\mathrm{d}}
\begin{document}
\section*{UnsolvedMath ID: 30003957}
\textbf{Source problem code:} catalogue entry
\subsection{Definitions and mathematical statement}
A closed manifold is a compact manifold without boundary. In the usual homogeneous-space setup, a geometric model consists of a manifold \(X\) and a Lie group \(G\) acting transitively on \(X\). A \((G,X)\)-structure on a manifold \(M\) of the same dimension as \(X\) is a maximal atlas of charts into \(X\) whose transition maps are locally restrictions of elements of \(G\). The supplied record does not specify the admissible geometric models, regularity, or further hypotheses on the action.

For connected \(M\), the usual developing-map description associates to such a structure a local diffeomorphism \(\operatorname{dev}\colon\widetilde M\to X\) and a holonomy homomorphism \(\rho\colon\pi_1(M)\to G\) satisfying
\[
\operatorname{dev}(\gamma\cdot p)=\rho(\gamma)\cdot\operatorname{dev}(p)
\qquad(\gamma\in\pi_1(M),\ p\in\widetilde M).
\]
The holonomy homomorphism is defined up to conjugation when no additional framing is chosen. The source does not specify whether \(M\) must be connected; the description above applies to a connected component.

Write \(\mathcal D_{G,X}(M)\) for the deformation space intended by the source and \(\operatorname{Rep}(\pi_1(M),G)\) for its intended representation space. A common marked deformation-space convention identifies structures under pullback by diffeomorphisms isotopic to the identity. Possible target conventions include \(\operatorname{Hom}(\pi_1(M),G)\), a conjugacy quotient, and, where applicable, an algebraic character quotient. These are different choices; the supplied record does not identify the intended equivalence relations or topologies. In particular, a map to homomorphisms themselves requires appropriate framing or normalization data.

\textbf{Source question, with unresolved conventions retained.} For every closed manifold \(M\) and every geometric model \((G,X)\) within the intended scope, consider the holonomy map
\[
\operatorname{Hol}\colon\mathcal D_{G,X}(M)\longrightarrow
\operatorname{Rep}(\pi_1(M),G).
\]
The source describes this map as open with discrete fibers. Treating these as supplied properties in the intended setting, is \(\operatorname{Hol}\) always a branched local homeomorphism? Find conditions under which it is a local homeomorphism.

Here openness means that \(\operatorname{Hol}(U)\) is open in the target for every open subset \(U\) of the domain. Discrete fibers means that each fiber, with its subspace topology, is discrete. Being a local homeomorphism at \(s\in\mathcal D_{G,X}(M)\) means that there are open neighborhoods \(U\) of \(s\) and \(V\) of \(\operatorname{Hol}(s)\) such that
\[
\operatorname{Hol}|_U\colon U\xrightarrow{\ \cong\ }V
\]
is a homeomorphism. The second request asks for hypotheses ensuring this property at specified structures or throughout the deformation space.

The source does not define ``branched local homeomorphism.'' A precise version of the first question therefore requires a specified branching definition, including any local-model, finiteness, or exceptional-set requirements. No such requirements are inferred here. The source's assertion of openness and discrete fibers also remains qualified by its unspecified deformation-space and representation-space conventions.
\subsection{Short English statement}
For geometric structures on a closed manifold, the source asks whether the holonomy map, described there as open with discrete fibers, always has a branched local homeomorphism structure, and which hypotheses make it a local homeomorphism. The intended spaces, geometric hypotheses, and meaning of branching require clarification before the first question has a unique precise formulation.
\subsection{Sources}
\begin{itemize}
\item Local UnsolvedMath registry, permanent ID \(30003957\), exact source problem code \texttt{catalogue entry}, title \emph{Local Branching of Holonomy Maps for Geometric Structures}; collection: Oberwolfach Reports Open Problems. The supplied record labels the problem \texttt{open} and the statement \texttt{corrected\_verified}, with a statement-verification note asserting preservation of mathematical content and repair of notation, typography, and omitted definitions. These are supplied registry assessments, not independent verification.
\item Supplied report citation: \emph{New Trends in Teichmüller Theory and Mapping Class Groups} (2019), Oberwolfach Reports, DOI \texttt{10.4171/owr/2018/40}, pp.\ 2475--2534. \url{https://doi.org/10.4171/owr/2018/40}. The background describes these reports as records of mathematical workshops and their open-problem discussions. The literature-source label is ``Oberwolfach Report 40, Alessandrini problem''; no proposer is specified in the proposer field. The supplied proposed year is 2018.
\item Supplied literature source: Goldman, \emph{geometric structures on manifolds}. \url{https://math.umd.edu/~wmg/gstom.pdf}.
\item Supplied literature source: \emph{Projective structures and quasi-Hitchin holonomy}. \url{https://doi.org/10.48550/arXiv.2110.15407}.
\item Supplied literature assessment dated exactly \texttt{2026-08-21}: the Ehresmann--Thurston principle gives local homeomorphism in standard compact settings, and special loci are understood, while the universal branched-local-homeomorphism claim and general criteria are described as unresolved. This is supplied dated curation derived from the OWR material, not independently verified. The source expressly requests independent verification of time-sensitive status and cited claims.
\end{itemize}
\end{document}
